% ===================================================================
% CHAPTER 9: LIMITATIONS AND FUTURE WORK
% ===================================================================
\chapter{Limitations and Future Work}
\label{ch:limitations}

\section{Limitations}

\subsection{Planar Fields}
% Ported from: cwaight_systems_paper.tex Sec. VI-C (Limitations);
% thesis.tex Ch.4 Limitations (lines 2348-2350); idetc2025_testbed.txt
% Sec. 5 (System Limitations, 2D constraint).

The zeroth-order primitives of Chapter~\ref{ch:zeroth}, hardware-tested
on the HSV-encoded testbed of Chapter~\ref{ch:tools}, and the
first-order attraction and orbital primitives of
Chapter~\ref{ch:first_order} and the second-order trackers of
Chapter~\ref{ch:second_order} are all restricted to planar vector
fields. Extending the first-order estimator to three dimensions would
require at minimum four non-coplanar robots, since an affine model in
three dimensions carries four unknowns per velocity component rather
than three. Extending the testbed's own HSV floor-map encoding to a
third dimension would separately require additional sensing capability
beyond the color sensor used throughout \cite{39}. No
three-dimensional extension of either the estimator or the testbed
representation is developed in this dissertation.

\subsection{No Advection}
\label{sec:ch9:no_advection}
% Ported from: cwaight_systems_paper.tex Sec. III-A (lines 213-215);
% Draft_11.tex Sec. VII (lines 1088-1094).

In both control architectures, the sensed field enters only the
estimation layer. It does not appear as a force or drift term in the
robot dynamics, so each robot's velocity responds to its commanded
value alone, and no physical flow acts on the platform. Every result in
Chapters~\ref{ch:first_order} and~\ref{ch:second_order}, including the
second-order tracker's run on real Santa Barbara Channel current data,
assumes the field is sensed rather than felt. That premise is exact for
a Decabot deployment, where the field is encoded and read rather than
physically experienced \cite{39}, but on a genuine ocean deployment
it is an extrapolation, since the current would advect the vehicle
directly, and it is untested against real advection in either chapter.

\subsection{Centralized Control and Team Size}
% Ported from: Draft_11.tex Sec. VII (lines 1061-1065);
% cwaight_systems_paper.tex Sec. VI-C (lines 484-485).

Cluster space control is centralized, requires global state, becomes
intractable for suitably large clusters, and is impractical without
reliable high-bandwidth communication among the robots \cite{53}.
Every result in this dissertation, three robots for the zeroth-order
and first-order primitives of Chapters~\ref{ch:zeroth}
and~\ref{ch:first_order}, and six for the
second-order trackers of Chapter~\ref{ch:second_order}, is a small
cluster in a limited workspace with ample bandwidth. Team size was not
varied within a fit order to probe this boundary; the estimation
framework of Chapter~\ref{ch:estimation} supports larger, overdetermined
teams at each order, but larger formations raise communication
bandwidth, computation, and formation-control-stability costs this
dissertation does not quantify. All three chapters also share a
centralized architecture; extending any of the three primitive families
to a decentralized, swarm-style architecture remains open, and the
testbed of Chapter~\ref{ch:tools} was not built with the controller
substitution such a comparison would need \cite{39}.

\subsection{Static and Simulated Fields}
% Ported from: cwaight_systems_paper.tex Sec. VI-C (lines 481-483);
% Draft_11.tex Sec. VII (line 1088); idetc2025_testbed.txt Sec. 5 (System
% Limitations, static-field restriction).

The testbed of Chapter~\ref{ch:tools}, and with it the zeroth-order
primitives of Chapter~\ref{ch:zeroth} and the first-order primitives of
Chapter~\ref{ch:first_order}, is restricted to static printed vector
fields; there is no clear path from a printed colormap to a reactive,
dynamic environment, so no hardware result in either chapter tests
performance against a field that changes during a trial
\cite{39}. Ground truth ties directly to the print, but the
trials do not exercise dynamic field behavior. The second-order trackers
of Chapter~\ref{ch:second_order} are evaluated on the steady double
gyre and on real, time-varying Santa Barbara Channel current data, but
every second-order result is simulated: no hardware trial of either
tracker exists, and the ocean run itself samples the gridded current
product without added measurement noise, so the noise-robustness
results of that chapter rest on the double gyre alone. More generally,
an indoor, small-scale testbed of this kind does not necessarily
represent the scale of a real oceanic or atmospheric deployment
\cite{39}.

\subsection{Sampling Rate and Communication}
% Ported from: cwaight_systems_paper.tex Sec. VI-C (lines 487-488);
% Draft_11.tex Sec. VII (centralized-communication clause,
% lines 1061-1065).

Both control architectures assume a sampling rate high enough that the
closed loop approximates continuous control, and both assume every
robot reports its position and field measurement synchronously each
cycle. Chapter~\ref{ch:first_order}'s hardware trials run this loop at
10~Hz over WiFi. Robustness to lower update rates, to communication
dropouts, to partial measurement sets, or to asynchronous reporting was
not tested in either chapter. Every robot also samples only a single
point in the field at its current location, with no look-ahead; this
limits testing of any predictive or path-planning extension to the
purely reactive primitives studied here \cite{39}.

\subsection{Physical Testbed Constraints}
% Ported from: idetc2025_testbed.txt Sec. 5 (System Limitations).

The Decabot testbed of Chapter~\ref{ch:tools} that hosts every hardware
result in this dissertation carries its own physical constraints,
independent of estimation order. The workspace is bounded between
roughly 1~m~$\times$~1~m and 3~m~$\times$~6~m, battery life limits a
continuous trial to about one hour per robot, and the commandable
velocity range is 0.05 to 0.3~m/s, set by the stiction floor and the
motor's rated maximum \cite{39}. The fleet is homogeneous and
carries no obstacle-avoidance sensing, so heterogeneous multi-robot
behavior and obstacle-rich environments are outside every hardware
result reported here. The printed floor map itself wears under repeated
robot traffic and needs periodic replacement, an operational cost
distinct from any algorithmic limitation \cite{39}.

\subsection{Estimator Conditioning and Tracker Robustness Near
Degeneracies}
\label{sec:ch9:estimator_conditioning}
% Ported from: Draft_11.tex Sec. VII (lines 1067-1085);
% cwaight_systems_paper.tex Sec. VI-A (condition number discussion,
% line 458).

Chapter~\ref{ch:estimation}'s formation matrix $\boldsymbol{\Phi}$ is
singular when the robots lie on a common curve of the fit's own degree,
a line for the first-order estimator and a conic for the second-order
estimator. The first-order estimator's degradation as the three-robot
formation approaches collinearity is discussed in
Chapter~\ref{ch:first_order}; monitoring $\kappa(\boldsymbol{\Phi})$
online and correcting for it is not implemented in either chapter. No
field tested in Chapter~\ref{ch:second_order} drives the six-robot
formation toward a common conic, so the second-order estimator is
untested in the regime its own conditioning limit describes.

Two further robustness gaps are specific to the second-order trackers.
The eigenvector comparison the $s_1$ tracker uses for tangent selection
weakens near saddle points; the capture latch of
Chapter~\ref{ch:second_order} damps, but does not exclude, the residual
risk of a wrong branch under heavy noise. The $D$ tracker's traversal
past a saddle is untested outside the double-gyre benchmark, and on the
benchmark it already runs outside that tracker's own stated premise,
since the fitted Hessian is locally indefinite where the true Hessian
is not.

\subsection{Second-Order Hardware Verification}
\label{sec:ch9:second_order_hw}
% Ported from: Draft_11.tex Sec. VII ("All results are simulated,"
% line 1088); thesis.tex Ch.5 Limitations (lines 4637-4638); parallel
% noted per Hart's Sec. 7.2 preliminary primitives.

Every second-order result in Chapter~\ref{ch:second_order}, on the
double gyre and on the Santa Barbara Channel record alike, is
simulation only. No hardware trial of the $D$ tracker or the $s_1$
tracker exists on the Decabot testbed of Chapter~\ref{ch:tools}, and
the calibration and dynamics identification that testbed provides have
so far validated only the zeroth- and first-order primitives of
Chapters~\ref{ch:zeroth} and~\ref{ch:first_order}. The second-order
trackers are, in this sense, in the same position the ridge, trench,
and saddle-point primitives of an earlier reactive-swarm architecture
occupied: developed and verified in simulation, with hardware
validation still to come \cite{59}.

\section{Future Work}

\subsection{Third-Order Estimation from Ten Robots}
% Ported from: Draft_11.tex Sec. VII, item c) (lines 1096-1100),
% generalizing the minimality argument of Chapter~\ref{ch:estimation}
% (Draft_11.tex Sec. II-B, lines 208-233).

The estimator of Chapter~\ref{ch:estimation} generalizes to a cubic
local model by the same counting argument that fixes the first- and
second-order minimums. A planar cubic basis has ten terms per velocity
component, $1, x, y, x^2, xy, y^2, x^3, x^2y, xy^2, y^3$, so ten robots
is the minimum formation size, and the formation matrix is singular
exactly when the ten sampling points lie on a common cubic curve, the
third-order analogue of the collinear and common-conic degeneracies of
Chapters~\ref{ch:first_order} and~\ref{ch:second_order}. A cubic fit
recovers the third derivatives of the field that the quadratic model of
Chapter~\ref{ch:second_order} truncates away. Those third derivatives
are exactly what is missing from the true Hessian of $D$ and from the
transverse curvature of an $s_1$ trench; recovering them would remove
the structural bias in $\hat{\mathbf{H}}_D$ noted in
Section~\ref{sec:ch9:ow_corner_branch} below, rather than only bound
it.

\subsection{Online Formation Reshaping}
% Ported from: Draft_11.tex Sec. VII, item b) (lines 1096-1098);
% thesis.tex Ch.5 Future Work (lines 4665-4668); cwaight_systems_paper.tex
% Sec. VI-A (fallback strategies, line 458).

Adapting a formation's shape online, the pair lengths and SAS
parameters of Chapter~\ref{ch:estimation}'s formations, to preserve the
conditioning of $\boldsymbol{\Phi}$ as a cluster nears a degenerate
configuration, is future work at every fit order. For the first-order
triangle of Chapter~\ref{ch:first_order}, this means resizing or
reshaping as the formation approaches collinearity; for the
second-order pentagon of Chapter~\ref{ch:second_order}, it means the
same as the formation approaches a common conic. Neither chapter
implements the fix; both leave it as the direct answer to the
conditioning gap of Section~\ref{sec:ch9:estimator_conditioning} above.

\subsection{Advection-Aware Control}
\label{sec:ch9:advection_control}
% Ported from: cwaight_systems_paper.tex Sec. VI-D (Future Work,
% line 496), paired with the No-Advection limitation above
% (Draft_11.tex Sec. VII, lines 1091-1094).

Extending the robot dynamics model of Chapter~\ref{ch:tools} to include
a direct advective term, so that a physical current displaces the
platform rather than only the sensed field, is a priority for future
work in both estimation chapters. This closes the gap identified in
Section~\ref{sec:ch9:no_advection} above: it is the difference between
the Decabot testbed, where no physical flow acts on the robot, and a
genuine ocean or atmospheric deployment, where it does.

A complementary route attacks the same sim-to-reality gap from the
robot side rather than the field side: a digital twin of the Decabot
platform, built by capturing each robot's velocity and acceleration
response directly rather than assuming the first-order momentum model
of Chapter~\ref{ch:tools}, would let the simulator used throughout
Chapters~\ref{ch:first_order} and~\ref{ch:second_order} reproduce
hardware dynamics more faithfully, independent of whether an advective
term is also added \cite{39}.

\subsection{Three-Dimensional Fields}
% Ported from: cwaight_systems_paper.tex Sec. VI-C. The second-order
% count and degeneracy condition are new text, the three-variable form of
% Chapter~\ref{ch:estimation}'s minimality and degeneracy results.
%% NOTE(new-math, checked 2026-09-29): not in any source paper.

A polynomial of total degree $k$ in three variables has
$\binom{k+3}{3} = (k+1)(k+2)(k+3)/6$ coefficients per field component,
so the minimality argument of Chapter~\ref{ch:estimation} carries over
with 4, 10, and 20 unknowns per component for $k = 1, 2, 3$. A
first-order estimate in three dimensions therefore needs at least four
robots, and it is degenerate exactly when they are coplanar. A
second-order estimate needs at least ten robots; with exactly ten, the
formation matrix is singular exactly when the robots lie on a common
quadric surface, since a null vector of the matrix is the coefficient
vector of a quadric through all ten positions, the three-dimensional
counterpart of the common-conic condition. No formation avoiding this
set has yet been designed, and the second-order surrogates themselves
change in three dimensions, where the Okubo-Weiss partition generalizes
to the $Q$-criterion and the strain tensor has three eigenvalues. The
zeroth-order primitives of Chapter~\ref{ch:zeroth} were also developed
and tested only in two dimensions.

\subsection{Additional Topological Features}
% Ported from: cwaight_systems_paper.tex Sec. VI-D (line 498);
% thesis.tex Ch.4 Future Work (lines 2375-2386).

Chapter~\ref{ch:second_order} tracks two Eulerian surrogates for a
separating boundary, the Jacobian determinant and the rate-of-strain
eigenvalue. Chapter~\ref{ch:preliminaries} introduces further
candidates, the Q-criterion and the Hua-Klein criterion among them,
that neither tracker in this dissertation follows. Heteroclinic orbits
connecting distinct critical points through their manifolds, offering
transit routes through a complex field, are also untouched; the
eigenvector-alignment argument Chapter~\ref{ch:first_order} sketches
for detecting such connections between neighboring critical points was
never implemented or tested. More generally, the feature taxonomy of
Chapters~\ref{ch:zeroth} through~\ref{ch:second_order} is not
exhaustive, and other regions of interest in a vector field remain to
be catalogued.

\subsection{Deployment on Surface Vessels}
% Ported from: Draft_11.tex Sec. VII, item d) (lines 1096-1100),
% paired with the Second-Order Hardware Verification limitation above;
% idetc2025_testbed.txt Sec. 6 (Future Work, outdoor surface-vessel and
% UAV testing).

Testing both second-order trackers on the Decabot testbed of
Chapter~\ref{ch:tools}, and then on surface vessels, is the direct
extension the simulation-only status of
Section~\ref{sec:ch9:second_order_hw} above calls for. A Decabot
deployment would use a field emitted by an onboard computer and sampled
by the robots' own sensors, following the pattern
Chapter~\ref{ch:tools} already establishes for the zeroth- and
first-order primitives. A surface-vessel deployment would additionally
confront the advection gap of
Section~\ref{sec:ch9:advection_control} directly, since a vessel's own
motion is no longer decoupled from the current it senses. Outdoor
testing was already proposed independently, for the zeroth- and
first-order primitives, as a way to examine transferability from the
indoor testbed's small scale to a real deployment, and it named both
autonomous surface vessels and UAVs as candidate platforms
\cite{39}; the second-order trackers would follow the same path.

\subsection{Scalar-Field Calibration and a Unified Testbed}
% Ported from: idetc2025_testbed.txt Sec. 6 (Future Work, greyscale
% sensor calibration for scalar fields).

The neural-network calibration technique of Chapter~\ref{ch:tools} was
developed for the hue-and-saturation encoding used throughout this
dissertation's vector-field trials. Extending the same supervised
calibration approach to greyscale sensors would support monochrome
colormaps encoding a scalar field, letting a single testbed and a
single calibration methodology serve both the vector-field primitives
of Chapters~\ref{ch:zeroth} through~\ref{ch:second_order} and the
scalar-field extremum-, contour-, and front-following primitives that
motivate them \cite{39}. No scalar-field calibration or hardware
trial is reported in this dissertation.

\subsection{Okubo-Weiss Corner Branch Selection (Deferred 2026-07-02)}
\label{sec:ch9:ow_corner_branch}

% Research note carried over unchanged from phd_thesis_template.tex's
% prior Chapter 5 (git history, this file). Relocated here because it
% is a specific, evaluable future-work item, not a general limitation.
% Sources: experiments/ow_clean_runs.py and
% experiments/outputs/ow_clean/ under
% trunk/Python_Simulations/Vector_Fields/VF_Robot/, git commits 7dc01a8
% and c19732f, and plan.md (repo root) Phase 3 and Phase 7 entries.

\textbf{Context.} The Okubo-Weiss boundary tracker of the second paper
drives $\hat{D} = \det(\hat{\mathbf{J}})$ to zero and drifts along the
level-set tangent $\mathrm{perp}(\nabla\hat{D})$, sign-aligned with the
ambient flow. On the steady double gyre the $D = 0$ set is not a pair of
closed diamonds but a network of straight lines,
$x_f \pm y_f \in \tfrac{1}{2} + \mathbb{Z}$, which cross at the diamond
corners where $\nabla D = \mathbf{0}$. Noise-free experiments (four
starts) showed lock-on within 13 to 16 control steps and mean $|D|$
between 0.003 and 0.007 on the boundary, 0.3 to 0.8 percent of the well
depth $\pi^4 A^2 = 0.974$, but the tangent follows its line straight
through every corner crossing: the gradient-floor fallback carries the
team across the degenerate point and the Newton action resumes on the
continuing line. Circulation around a single vortex core therefore never
occurs with the published law. The paper presents this straight-through
behavior honestly; the wording of Problem 2 was changed from ``circulate
along'' to ``travel along'' the boundary.

\textbf{The deferred idea.} At an X-crossing the field is locally
$D(\mathbf{p}^* + \mathbf{t}) \approx \tfrac{1}{2}\mathbf{t}^\top
\mathbf{H}_D \mathbf{t}$ with $\mathbf{H}_D$ indefinite, so the zero set
consists of the two asymptote directions $\mathbf{t}$ satisfying
$\mathbf{t}^\top \mathbf{H}_D \mathbf{t} = 0$. In the eigenbasis of
$\mathbf{H}_D$ with $\lambda_1 < 0 < \lambda_2$ these lie at angles
$\pm\arctan\sqrt{-\lambda_1/\lambda_2}$ from the $\lambda_2$
eigenvector. A corner rule can therefore be computed onboard from
$\hat{\mathbf{H}}_D$ alone: inside the gradient-floor ball, evaluate
both estimated asymptotes, discard the one parallel to the incoming
tangent, and exit along the turning branch, sign-chosen by the flow.
Selecting the turning branch at every corner closes the circuit around
one diamond; always selecting the continuing branch reproduces the
published straight-through behavior, so the rule is a second
one-parameter behavior switch, mirroring the park-versus-traverse knob
of the separatrix tracker.

\textbf{Known risk.} The $\hat{\mathbf{H}}_D$ estimate carries a
formation-radius-independent structural bias (third-derivative terms
absent from the quadratic model; see the estimator accuracy results of
the second paper). Eigen-directions survive that bias to within a few
tenths of a degree on the separatrix, but the accuracy of the estimated
asymptote directions at the corners specifically is untested, and near
the corners the eigenvalues themselves are strongly biased. This is the
first thing to measure.

\textbf{Validation plan when revisited.} Implement as a new primitive
(working name \textit{logic\_g\_fable\_step}) alongside the untouched
\textit{logic\_g\_newton\_contour\_pentagon}, following the pattern of
\textit{separatrix\_logic\_fable\_step} (bit-exact equivalence to the
parent law when the knob is off, then the new branch enabled).
Head-to-head at all eight diamond corners plus the domain-center
crossing: success is loop closure around one diamond; metrics are
circumnavigation time and mean $|D|$ per circuit, against the
straight-through baseline of \textit{ow\_clean\_runs.py}.
